\documentclass[11pt]{article}

% --- Encoding & fonts ---
\usepackage[utf8]{inputenc}
\usepackage[T1]{fontenc}

% --- Japanese support (compile with lualatex/xelatex, or use CJK with pdflatex) ---
% If you compile with pdflatex, CJKutf8 lets you typeset Japanese inline.
\usepackage{CJKutf8}

% --- Layout & math ---
\usepackage[margin=2.5cm]{geometry}
\usepackage{amsmath, amssymb}
\usepackage{booktabs}
\usepackage{enumitem}
\usepackage{hyperref}
\hypersetup{
    colorlinks=true,
    linkcolor=blue,
    citecolor=blue,
    urlcolor=blue,
}

% --- Bibliography ---
\usepackage[backend=biber, style=numeric]{biblatex}
\addbibresource{references.bib}

\title{Japanese Notes --- JLPT N5}
\author{}
\date{\today}

\begin{document}
\begin{CJK}{UTF8}{min}

\maketitle
\tableofcontents

\section{Introduction}
JLPT N5 is the lowest level of the Japanese-Language Proficiency Test. These
notes cover the core hiragana, katakana, basic vocabulary, and grammar needed
to reach this level.

\section{Hiragana}
% Example of typesetting Japanese: あ い う え お

\section{Katakana}

\section{Vocabulary}

\section{Grammar}

\section{Diary}



% --- Bibliography ---
\printbibliography

\end{CJK}
\end{document}
